% Part: lambda-calculus
% Chapter: church-rosser
% Section: beta-eta-reduction

\documentclass[../../../include/open-logic-section]{subfiles}

\begin{document}

\olfileid{lam}{cr}{be}

\olsection{Reduksi-$\beta\eta$}

Sipat Church--Rosser
uga berlaku kanggo
reduksi-$\beta\eta$ ($\bered$).

\begin{lem}\ollabel{lem:one-par}
  Yen $M \beredone M'$,
  mula $M \beredpar M'$.
\end{lem}

\begin{proof}
  Gunakake induksi ing !!{derivation}
  saka $M \beredone M'$.
  Yen $M \eredone M'$
  lumantar konversi-$\eta$
  (yaiku
  \olref[syn][eta]{defn:beredone}),
  gunakake
  \olref[pbe]{thm:refl}
  kanggo badan suku banjur
  aturan kontraksi eta paralel.
  Kontraksi ing konteks kaangkat
  dening aturan kompatibilitas.
  Kasus liyane kaya ing
  \olref[b]{lem:one-par}.
\end{proof}

\begin{lem}\ollabel{lem:par-red}
  Yen $M \beredpar M'$,
  mula $M \bered M'$.
\end{lem}

\begin{proof} Gunakake induksi ing
  !!{derivation} saka
  $M \beredpar M'$.

  Yen aturan pungkasan
  \olref[pbe]{defn:beredpar5},
  $M$ awujud
  $\lambd[x][Nx]$
  lan $M'$ awujud
  $N'$ kanggo sawetara
  $x$, $N$, $N'$ kanthi
  $x \notin FV(N)$ lan
  $N \beredpar N'$.
  Mula luwih dhisik
  $\lambd[x][Nx]$
  direduksi dadi $N$
  lumantar konversi-$\eta$,
  banjur diterusake
  urutan langkah
  $\beredone$ kang
  nuduhake $N \bered N'$,
  kaya kang diwenehake
  hipotesis induksi.
  Patang kasus liyane
  ngetutake pola simulasi
  reduksi paralel sadurunge.
\end{proof}

\begin{lem}\ollabel{lem:str}
  $\bered$ yaiku relasi
  transitif paling cilik
  kang ngemot $\beredpar$.
\end{lem}

\begin{proof}
  Kaya ing \olref[b]{lem:str}.
\end{proof}

\begin{thm}\ollabel{thm:cr}
  $\bered$ nduweni
  sipat Church--Rosser.
\end{thm}

\begin{proof}
  Miturut \olref[dap]{thm:str},
  \olref[pbe]{thm:cr},
  lan \olref{lem:str}.
\end{proof}
\end{document}
